\documentclass[a4paper,12pt]{article} 
\usepackage[utf8]{inputenc} % Acentos válidos sin problemas
\usepackage[spanish]{babel} % Idioma
%-----------------------------------INICIO DE PACKETES------------------/
%----------------------------------------------------------------------/|
\usepackage{amsmath}   % Matemáticas: Comandos extras(cajas ecuaciones) |
\usepackage{amssymb}   % Matemáticas: Símbolos matemáticos              |
\usepackage{datetime}  % Agregar fechas                                 |
\usepackage{graphicx}  % Insertar Imágenes                              |
%\usepackage{biblatex} % Bibliografía                                   |
\usepackage{multicol}  % Creación de tablas                             |
\usepackage{longtable} % Tablas más largas                              |
\usepackage{xcolor}    % Permite cambiar colores del texto              |
\usepackage{tcolorbox} % Cajas de color                                 |
\usepackage{setspace}  % Usar espacios                                  |
\usepackage{fancyhdr}  % Para agregar encabezado y pie de página        |
\usepackage{float}     % Flotantes                                      |
\usepackage{soul}      % "Efectos" en palabras                          |
\usepackage{hyperref}  % Para usar hipervínculos                        |
\usepackage{caption}   % Utilizar las referencias                       |
\usepackage{subcaption} % Poder usar subfiguras                         |
\usepackage{multirow}  % Nos permite modificar tablas                   |
\usepackage{array}     % Permite utilizar los valores para multicolumn  |
\usepackage{booktabs}  % Permite modificar tablas                       |
\usepackage{diagbox}   % Diagonales para las tablas                     |
\usepackage{colortbl}  % Color para tablas                              |
\usepackage{listings}  % Escribir código                                |
\usepackage{mathtools} % SIMBOLO :=                                     |
\usepackage{enumitem}  % Modificar items de Listas                      |
\usepackage{tikz}      %                                                |
\usepackage{lipsum}    % for auto generating text                       |
\usepackage{afterpage} % for "\afterpage"s                              |
\usepackage{pagecolor} % With option pagecolor={somecolor or none}|     | 
%\usepackage{glossaries} %                                              |        
%----------------------------------------------------------------------\|
%-----------------------------------FIN--- DE PACKETES------------------\
%\usepackage{newtxtext,newtxmath}
%\usepackage{bm}
%--------------------------------/
%-------------------------------/
\usepackage[                 %   |
  headheight=15pt,  %            |
  letterpaper,  % Tipo de pag.   |
  left =1.5cm,  %  < 1 >         |
  right =1.5cm, %  < 1 >         | MARGENES DE LA PAGINA
  top =2cm,     %  < 1.5 >       |
  bottom =1.5cm %  < 1.5 >       |
]{geometry}     %                |
%-------------------------------\
%--------------------------------\

%----------------------------------------------------------------------/
%-------------------Encabezado y Pie de Pagina-----------------------/ |
%--------------------------------------------------------------------\ |
\usepackage{fancyhdr} %                                                |
\usepackage{lastpage} %                                                |
%\fancyhf{}           %                                                |
\pagestyle{fancy}     %                                                |
\fancyhead[L]{{$2024-1$}} %                                              |
\fancyhead[C]{}                      
\fancyhead[R]{}

\fancyfoot[L]{} 
\fancyfoot[C]{\thepage}
\fancyfoot[R]{} 

\renewcommand{\headrulewidth}{1.5pt} %                                 |
\renewcommand{\footrulewidth}{1.5pt} %                                 | 
%--------------------------------------------------------------------\ |
%-------------------Encabezado y Pie de Pagina-----------------------/ |
%----------------------------------------------------------------------/


\begin{document}%----------------------INICIO DOCUMENTO------------|
%------------------------------------------------------------------|

\newpage

\begin{titlepage}
    \center% Center everything on the page
    % Defines a new command for the horizontal lines, change thickness here
    \newcommand{\HRule}{\rule{\linewidth}{0.5mm}} 
    
    %---------------------
    %	LOGO
    %---------------------
    \includegraphics[width=4.5cm]{IMA/Ciencias.png} \\ 
    
    %----------------------------
    %	Encabezado
    %----------------------------
    \quad \\[1.5cm]
    %\textsc{\LARGE MSc Thesis}\\[1.5cm] % Name of your university/college
    \textsc{\Large Facultad de Ciencias}\\[0.5cm] % Major heading such as course name
    \textsc{\Large Computación Distribuida}\\[0.5cm] % Major heading such as course name
    %\textsc{\large Del Aula a la START UP}\\[0.5cm] % Minor heading such as course title
    %-------------
    %	TITULOS
    %-------------
    \makeatletter
        \HRule\\ [0.4cm]
            { \huge \bfseries PRACTICA 02}\\
        \HRule\\ [0.4cm]
    %-------------------
    %	PROFESORES
    %-------------------
    %\vspace{10mm}
    \large\textbf{Semestre $2024-1$}\\
    
    \vspace{10mm}
    \begin{minipage}{0.7\textwidth}
        \begin{flushleft} \large
            \emph{Profesor:}\\
                Luis Germán Pérez Hernández\\
                \vspace{3mm}
                \emph{Ayudantes:}\\
                Daniel Michel Tavera\\
                Yael Antonio Calzada Martín\\
        \end{flushleft}
    \end{minipage}
    %~
    \\
    %-------------------
    %	AUTOR
    %-------------------
    \vspace{10mm}
    \begin{minipage}{0.4\textwidth}
        \begin{flushright} \large
        \centering
            \emph{Alumnos} \\ % \emph{Alumnos:} \\
            \vspace{3mm}
            Marco Silva Huerta\\
            316205326\\
            \vspace{4mm}
            Edgar Montiel Ledesma\\
            317317794\\
            \vspace{4mm}
            Carlos Cortés\\
            420004846
            %\vspace{4mm}
            % Autor 2\\
        \end{flushright}
    \end{minipage}\\[1.2cm]
    \makeatother
    %---------------
    %	F E C H A
    %---------------
    {\large 18 de Octubre de 2023}\\[2cm]
    \vfill 
    \end{titlepage}

\newpage

\begin{center}
    {\Large Algotitmo Dijkstra Distribuido}
\end{center}

\section*{Forma de compilar}
para compilar el archivo dentro del archivo Practica01
\begin{itemize}
   \item[ ] cp@linux:cd Practica01/src
   \item[ ] \$ mpicc nombre-de-la-practica.c -o nombre-de-la-compilacion
   \item[ ] \$ ./nombre-de-la-compilacion
\end{itemize}

\section*{Funcionamiento}

Este código implementa el algoritmo de Dijkstra para encontrar el camino más corto en un grafo ponderado con múltiples procesos utilizando MPI

\begin{itemize}
   \item[1.] Incluye las bibliotecas necesarias, incluyendo `<stdio.h>`, `<stdlib.h>`, `<limits.h>`, `<time.h>` y `<mpi.h>`.

   \item[2.] Define la constante `INFINITY` como `INT MAX`, que se utiliza para representar una distancia infinita en el algoritmo de Dijkstra.

   \item[3.] Define una función `randomDelay` que genera un valor de retraso aleatorio entre 1 y 1000.

   \item[4.] Define una función `minDelay` que encuentra el vértice con el retraso mínimo no visitado en el grafo.

   \item[5.] Define una función `printDijkstra` para imprimir los resultados del algoritmo de Dijkstra, incluyendo el retraso mínimo y la ruta desde el nodo 0 hasta cada nodo.

   \item[6.] La función `dijkstra` implementa el algoritmo de Dijkstra en paralelo. Cada proceso MPI calcula las distancias parciales desde el nodo 0 hasta un subconjunto de nodos. Luego, se sincronizan los resultados entre los procesos para obtener la distancia mínima global.

   \item[7.] La función `main` es la entrada principal del programa. Inicializa MPI, obtiene el rango y el tamaño del comunicador MPI, y solicita al usuario el número de nodos para el grafo.

   \item[8.] Utiliza MPI para transmitir el número de nodos a todos los procesos.

   \item[9.] Crea una matriz `graph` para representar el grafo y llena sus valores con retrasos aleatorios. Esto se hace solo en el proceso con rango 0.

   \item[10.] Utiliza MPI para transmitir la matriz `graph` a todos los procesos.

   \item[11.] Llama a la función `dijkstra` para calcular el camino más corto desde el nodo 0 hasta todos los demás nodos en paralelo.

   \item[12.] Libera la memoria asignada para la matriz `graph` y finaliza MPI.
\end{itemize}

\section*{Pseudocódigo del Algotitmo}
\begin{verbatim}
    1. Inicializar todas las distancias en D con un valor infinito relativo, 
       ya que son desconocidas al principio, exceptuando la de a, 
       qué se debe colocar en 0, pues la distancia de a a si mismo sería 0. 
       C es copia de V        
    
    2. Para todo vértice i en C se establece [PI]= a.        
    
    3. Se obtiene el vértice s en C tal que no existe otro vértice w en C 
       tal que (D[w] < D[s]). 
       Para esto se envía un mensaje al nodo correspondiente y se regresa
       un mensaje de respuesta en donde se toma el tiempo y se le asigna a 
       su distancia correspondiente. De manera concurrente el nodo destino 
       realiza el mismo procedimiento para calcular su distancia a sus nodos
       vecinos que no han sido visitados.
    
        // En lugar de buscar el vértice con la distancia más corta 
        // iterativamente, ahora se utiliza una heap para mantener una lista 
        // de vértices no visitados, ordenada por la distancia más corta. Así 
        // encontrar el vértice n la distancia más corta en tiempo logarítmico.    
    
    4. Se elimina de C el vértice s. El vértice u se elimina del conjunto C.
    
    5. Para cada arista e en E de longitud l, 
       que une el vértice s con algún otro vértice t en C,
       Para cada arista que sale del vértice u, se verifica si la distancia 
       a través del vértice u es menor que la distancia actual del vértice t.
    
        - Si l + D[s] < D[t], entonces:
         // Si la distancia a través del vértice u es menor que la distancia 
         // actual del vértice t, entonces se actualiza la distancia del vértice t.
        - Se establece D[t] :=  l + D[s].
         // La distancia del vértice t se establece en la suma de la distancia 
         // del vértice u y el peso de la arista.
        - Se establece P[t] := s.
        // El predecesor del vértice t se establece en el vértice u.
    
    6. Se regresa al paso 4.
        // El algoritmo regresa al paso 4 y repite el proceso hasta que todos
        // los vértices hayan sido visitados.    
\end{verbatim}
    

\section*{Desarrollo}








\end{document}%---------------------- FIN DOCUMENTO---------------|
%-----------------------------------------------------------------|